% VLDB WORKSHOP template version of 2024-05-15 enhances the ACM template, version 1.7.0:
% https://www.acm.org/publications/proceedings-template
% The ACM Latex guide provides further information about the ACM template.

\documentclass[sigconf, nonacm]{acmart}

\newcommand\vldbyear{2026}
\newcommand\vldbworkshop{Quality in Databases (QDB)}
\newcommand\vldbauthors{\authors}
\newcommand\vldbtitle{\shorttitle}
\newcommand\vldbavailabilityurl{}
\newcommand\vldbpagestyle{plain}

\begin{document}

\title[Credibility Measurement in Social Media]{Measuring Credibility in Social Media Platforms through Data Quality and Provenance}

\author{Sebasti\'an Garc\'ia Parra}
\affiliation{%
  \institution{Universidad de la Rep\'ublica}
  \city{Montevideo}
  \country{Uruguay}
}
\email{TODO@example.org}

\begin{abstract}
Social media platforms have become primary sources of information for health,
politics, crises, and everyday decisions. Yet users often consume and propagate
content without objective indicators about its credibility, especially when
content is copied, revised, or shared across platforms. This paper presents a
data-quality model for measuring the credibility of social media clips. The model
organizes credibility around two main dimensions: trustworthiness, computed from
reputation, verifiability, and expertise, and provenance, which captures the
origin and evolution of a clip. We introduce Trustworthiness Path Stability, a
quality metric that evaluates whether trustworthiness remains stable along a
provenance path. We also describe a modular processing flow that normalizes clips
from heterogeneous platforms, generates domain-specific features, computes
quality metrics, and reconstructs provenance using explicit interactions and
textual equivalence. A proof of concept on posts about statins and cholesterol,
validated with a medical-domain expert, shows high agreement between expert
judgments and the proposed model, while highlighting how provenance can change
the final credibility assessment.
\end{abstract}

\maketitle

%%% do not modify the following VLDB block %%
%%% VLDB block start %%%
\pagestyle{\vldbpagestyle}
\begingroup\small\noindent\raggedright\textbf{VLDB Workshop Reference Format:}\\
\vldbauthors. \vldbtitle. VLDB \vldbyear\ Workshop: \vldbworkshop.\\
\endgroup
\begingroup
\renewcommand\thefootnote{}\footnote{\noindent
This work is licensed under the Creative Commons BY-NC-ND 4.0 International License. Visit \url{https://creativecommons.org/licenses/by-nc-nd/4.0/} to view a copy of this license. For any use beyond those covered by this license, obtain permission by emailing \href{mailto:info@vldb.org}{info@vldb.org}. Copyright is held by the owner/author(s). Publication rights licensed to the VLDB Endowment. \\
\raggedright Proceedings of the VLDB Endowment.
ISSN 2150-8097. \\
}\addtocounter{footnote}{-1}\endgroup
%%% VLDB block end %%%

%%% do not modify the following VLDB block %%
%%% VLDB block start %%%
\ifdefempty{\vldbavailabilityurl}{}{
\vspace{.3cm}
\begingroup\small\noindent\raggedright\textbf{VLDB Workshop Artifact Availability:}\\
The source code, data, and/or other artifacts have been made available at \url{\vldbavailabilityurl}.
\endgroup
}
%%% VLDB block end %%%

\section{Introduction}

Social media content is increasingly used as evidence for decisions that have
online and offline consequences. Users may decide whether to share a post, trust
a source, accept medical advice, or reject a treatment based on information whose
quality is uncertain. This risk is amplified by the way social media platforms
encourage rapid sharing and by the lack of standards for preserving quality
metadata when content crosses platform boundaries.

Existing platform mechanisms, such as verified accounts or community annotations,
provide useful signals but do not provide a general data-quality view of
credibility. Moreover, users commonly participate in several platforms, and
content may be copied manually, quoted, paraphrased, or reposted without carrying
the metadata available in the original platform. Thus, the same information item
may appear in multiple versions, with different authors, evidence, audience
signals, and trust contexts.

This paper condenses a master's thesis on measuring credibility in social media
platforms. Its main contribution is a provenance-aware data-quality model for
clips, where a clip denotes a social media information unit such as a post, tweet,
comment, video caption, or platform-specific publication. The model is platform
agnostic and domain adaptable: a data-quality expert defines the dimensions and
metrics, while a domain expert adjusts weights and thresholds for a concrete use
case.

The paper makes three contributions. First, it defines a quality model for
credibility that combines trustworthiness and provenance. Second, it introduces a
metric, Trustworthiness Path Stability, to capture how trust changes as content is
propagated and transformed. Third, it presents a modular processing flow and a
proof of concept in the medical domain, focused on social media content about
statins and cholesterol.

\section{Background and Requirements}

Credibility has long been treated as a perceived quality rather than an inherent
property of an object~\cite{FoggTseng1999}. In data-quality terms, credibility is
related to trustworthiness, reputation, and provenance~\cite{BatiniScannapieco2016}.
For social media, this relation is especially important because credibility is
evaluated under uncertainty: users often see only a fragment of the original
context, and the perceived authority of an author may be misleading.

Prior work on information credibility in Twitter/X has shown that message
features, user features, topic features, and propagation features can help
estimate credibility~\cite{Castillo2011}. Other work studies social media data
quality directly and adapts traditional data-quality dimensions to streams of
tweets~\cite{Arolfo2020}. Recent user studies also show that source expertise and
prior beliefs strongly affect how readers judge social media posts, particularly
in health-related domains~\cite{Kuutila2024}.

Provenance provides a complementary view. Data provenance captures origin,
derivation, and transformations, and has been recognized as a central mechanism
for assessing trust in data products~\cite{Simmhan2005}. PROV-SAID extends W3C
PROV to represent information diffusion in social media~\cite{Taxidou2015}. It
models messages, agents, influence, and interactions, but its original formulation
is centered on information diffusion within social platforms and does not directly
model cross-platform sharing or quality metrics attached to the provenance path.

From these observations, we derive four design requirements. The model must be
platform agnostic, because credibility signals vary across X, Facebook, TikTok,
Reddit, WhatsApp, and other platforms. It must be provenance aware, because the
same content may become more or less credible as it is copied, revised, or shared.
It must support domain-specific quality metrics, because credibility in medicine,
politics, and emergency response depends on different evidence. Finally, it must
be modular, so new feature generators, metrics, and provenance search strategies
can be incorporated without redesigning the whole flow.

\section{Credibility Quality Model}

We model \emph{Credibility} as a quality cluster composed of two main dimensions:
\emph{Trustworthiness} and \emph{Provenance}. Trustworthiness estimates the
expected reliability of the clip in its domain. Provenance captures the origin,
derivation, and propagation path of the clip.

Trustworthiness is decomposed into three factors. \emph{Reputation} evaluates
external perception of the author or source, using metrics such as engagement and
confidence. \emph{Verifiability} evaluates whether the clip provides enough
evidence to corroborate its claims, including source data, modification
transparency, and verification metadata. \emph{Expertise} evaluates whether the
author has domain-specific knowledge, using metrics such as certification,
production, influence, and experience.

Provenance is decomposed into attribution, process, and trustworthiness path.
Attribution identifies the sources or entities that authored a clip. Process
captures how the clip was generated or modified. Trustworthiness path captures
how trustworthiness evolves across the provenance chain.

\begin{table}[t]
  \caption{Core dimensions and factors in the credibility model.}
  \label{tab:model}
  \begin{tabular}{lll}
    \toprule
    Dimension & Factor & Example metrics\\
    \midrule
    Trustworthiness & Reputation & Engagement, confidence\\
    Trustworthiness & Verifiability & Access to source data\\
    Trustworthiness & Expertise & Certification, influence\\
    Provenance & Attribution & Attribution metadata\\
    Provenance & Process & Edit count, edit history\\
    Provenance & Path & Path stability\\
    \bottomrule
  \end{tabular}
\end{table}

Let $T$ denote trustworthiness and $P$ denote provenance, both in $[0,1]$. The
credibility value $C$ is computed as a weighted combination:
\begin{equation}
  C = tT + pP,
\end{equation}
where $t$ and $p$ are domain-specific weights. Trustworthiness is similarly
computed from verifiability $V$, reputation $R$, and expertise $E$:
\begin{equation}
  T = vV + rR + eE.
\end{equation}
Weights can be set manually by a domain expert, inferred through rules, or learned
from expert-labeled clips.

\section{Provenance-Aware Processing Flow}

The processing flow receives a social media clip and computes its credibility
through six stages. First, the clip is ingested from a platform-specific source.
Second, it is normalized into a JSON representation that includes identifiers,
platform URI, content, dates, geolocation, interactions, media, comments,
performance metrics, user metadata, generated features, and quality metrics.
Third, feature modules extract domain-specific signals. Fourth, quality modules
compute metrics for each quality factor. Fifth, the provenance module searches for
equivalent clips and reconstructs the provenance graph. Finally, the credibility
module combines the metrics into a final score.

The provenance module distinguishes explicit and implicit interactions. Explicit
interactions are represented by platform mechanisms such as retweets, replies,
quotes, or shares. Implicit interactions occur when users copy, paraphrase, or
repost content without preserving a direct platform link. These cases are common
in cross-sharing, where content moves between platforms and may lose its original
metadata.

To address implicit interactions, the flow uses clip equivalence. For text clips,
we instantiate equivalence with TF-IDF cosine similarity~\cite{Yuntao2005}. Given
two clips $c_1$ and $c_2$, their textual similarity is:
\begin{equation}
  Sim(c_1,c_2) = \frac{c_1 \cdot c_2}{\|c_1\|\|c_2\|}.
\end{equation}
Equivalent clips are then mapped to a provenance representation inspired by
PROV-SAID. We adapt PROV-SAID by replacing messages with clips, adding a
\texttt{SocialNetwork} attribute, modeling cross-platform derivations, and
attaching quality metrics such as trustworthiness, verifiability, expertise, and
reputation to provenance entities.

\subsection{Trustworthiness Path Stability}

The main provenance-derived metric is \emph{Trustworthiness Path Stability}. It
measures whether trustworthiness remains stable as a clip propagates. Let
$T_i$ be the trustworthiness value of the $i$-th clip in a provenance path of
length $n$. We define:
\begin{equation}
  TPS = 1 - \frac{\sum_{i=1}^{n-1}|T_{i+1}-T_i|}{n-1}.
\end{equation}
Values close to 1 indicate stable trustworthiness across the path. Lower values
indicate fluctuations, possibly caused by content modifications, loss of source
quality, or propagation through less trustworthy actors.

\section{Proof of Concept}

We implemented a proof of concept for health-related social media clips about
statins and cholesterol. This domain was selected because misinformation about
statins can influence treatment adherence and public-health decisions. The
implementation was deployed on Google Cloud Platform using Pub/Sub for ingestion,
Cloud Storage for persistence, Cloud Functions for feature and quality modules,
and Cloud Composer for orchestration.

The experiment started from 97 clips collected from multiple social media
platforms using keywords related to statins and cholesterol. Eleven clips were
preselected to cover different platforms and credibility levels. A medical-domain
expert classified the selected clips into credibility levels and provided scores.

For the experiment, trustworthiness was computed from three metrics: access to
source data, certification, and confidence. The domain expert assigned weights of
0.4, 0.4, and 0.2, respectively. Credibility combined Trustworthiness Path
Stability and trustworthiness with weights of 0.6 and 0.4 when provenance was
available. When no provenance was available, the model relied only on
trustworthiness.

\begin{table}[t]
  \caption{Comparison between expert judgments and model output.}
  \label{tab:results}
  \begin{tabular}{ccccc}
    \toprule
    Clip & Expert & Score & Model & Cred.\\
    \midrule
    1 & Medium & 0.60 & High & 0.80\\
    2 & Medium & 0.70 & Medium & 0.68\\
    3 & Medium & 0.60 & Medium & 0.68\\
    4 & Medium & 0.70 & Medium & 0.64\\
    5 & Medium & 0.60 & High & 0.78\\
    6 & Low & 0.00 & Low & 0.09\\
    7 & Low & 0.00 & Low & 0.18\\
    8 & Low & 0.00 & Low & 0.11\\
    9 & Low & 0.10 & Low & 0.09\\
    10 & Low & 0.00 & Low & 0.08\\
    11 & Low & 0.00 & Low & 0.36\\
    \bottomrule
  \end{tabular}
\end{table}

The model agreed with the expert level in nine of the eleven clips
(\autoref{tab:results}). The two differences corresponded to clips for which
provenance increased the credibility score: both had stable paths and were shared
by users with relatively high trustworthiness. The expert did not have access to
all profiles involved in the cross-platform propagation path, which explains why
the expert score was closer to trustworthiness alone than to the final
provenance-aware score.

The preliminary implementation had an average execution time of 5220 ms per clip.
Although this result is encouraging, the sample is too small for claims about
scalability or predictive accuracy. The value of the experiment is instead to show
that a data-quality model can produce interpretable credibility metrics and expose
where provenance changes the assessment.

\section{Discussion}

The proposed model shifts credibility assessment from a monolithic prediction task
to a data-quality workflow. This has two advantages for QDB. First, it makes the
criteria explicit: users can inspect whether a score came from source
verifiability, author expertise, reputation, or provenance stability. Second, it
separates domain knowledge from infrastructure. Medical credibility can use
certification and access to source data, while other domains may define different
features and thresholds.

There are also limitations. The proof of concept uses a small set of clips and a
single domain expert. The clip-equivalence strategy is limited to text and does not
yet cover images or video. The provenance reconstruction depends on platform APIs,
search engines, and available public metadata. Finally, the model assumes that
domain experts can define meaningful weights; future work should compare manual,
rule-based, and machine-learned weight calibration.

Future work includes scaling the experiment, incorporating multimodal equivalence,
improving real-time processing, adding quality metadata for provenance completeness
and uncertainty, and designing user-facing visualizations for credibility metrics.

\section{Conclusion}

This paper presented a provenance-aware data-quality model for measuring
credibility in social media platforms. By combining trustworthiness, provenance,
clip equivalence, and Trustworthiness Path Stability, the model provides an
interpretable way to evaluate social media content across platforms and domains.
The proof of concept in the statins domain suggests that provenance can materially
change credibility assessment and should be treated as a first-class quality signal.

\begin{acks}
TODO: Add funding, advisor, institutional acknowledgements, and domain-expert
acknowledgements as appropriate.
\end{acks}

\bibliographystyle{ACM-Reference-Format}
\bibliography{references}

\end{document}
\endinput
